
\documentclass{article}
\usepackage{colortbl}
\usepackage{makecell}
\usepackage{multirow}
\usepackage{supertabular}

\begin{document}

\newcounter{utterance}

\centering \large Interaction Transcript for game `cladder', experiment `full\_v1.5\_default', episode 6302 with score{-}Qwen3.5{-}9B{-}sp.
\vspace{24pt}

{ \footnotesize  \setcounter{utterance}{1}
\setlength{\tabcolsep}{0pt}
\begin{supertabular}{c@{$\;$}|p{.15\linewidth}@{}p{.15\linewidth}p{.15\linewidth}p{.15\linewidth}p{.15\linewidth}p{.15\linewidth}}
    \# & \multicolumn{2}{c}{Player} && \multicolumn{2}{c}{Game Master} \\
    \hline

    \theutterance \stepcounter{utterance}  
    & & & \multicolumn{4}{p{0.6\linewidth}}{
        \cellcolor[rgb]{0.9,0.9,0.9}{
            \makecell[{{p{\linewidth}}}]{
                \texttt{\tiny{[P1$\langle$GM]}}
                \texttt{You are an expert in causal inference. The following question is not a typical commonsense query, but rather a meticulously designed question created by a professor specializing in causal inference, intended to assess the students' mastery of the course content.} \\
\texttt{Imagine a self{-}contained, hypothetical world with only the following conditions, and without any unmentioned factors or causal relationships: Zuph has a direct effect on rixq and xevu. Rixq has a direct effect on xevu. For those who are not zuph and are not rixq, the probability of xevu is 15\%. For those who are not zuph and are rixq, the probability of xevu is 50\%. For those who are zuph and are not rixq, the probability of xevu is 32\%. For those who are zuph and are rixq, the probability of xevu is 85\%. The overall probability of zuph is 84\%. Will rixq increase the chance of xevu?} \\
\texttt{Start your answer with 'Yes' or 'No', followed by additional reasoning or evidence to support your explanation.} \\
            }
        }
    }
     \\ \\

    \theutterance \stepcounter{utterance}  
    & \multicolumn{4}{p{0.6\linewidth}}{
        \cellcolor[rgb]{0.561,0.737,0.561}{
            \makecell[{{p{\linewidth}}}]{
                \texttt{\tiny{[P1$\rangle$GM]}}
                \texttt{Yes} \\
\\ 
\texttt{**Reasoning:**} \\
\\ 
\texttt{To determine if \$Rixq\$ increases the chance of \$Xevu\$, we must calculate the causal effect of \$Rixq\$ on \$Xevu\$. In this hypothetical world, \$Zuph\$ is a confounder because it affects both \$Rixq\$ (direct effect mentioned) and \$Xevu\$. Therefore, we cannot simply compare the marginal probabilities of \$Xevu\$ given \$Rixq\$; we must calculate the **Average Causal Effect (ACE)**, which is the difference in the expected probability of \$Xevu\$ under intervention (\$do(Rixq=1)\$) versus no intervention (\$do(Rixq=0)\$).} \\
\\ 
\texttt{The problem provides the conditional probabilities \$P(Xevu | Zuph, Rixq)\$:} \\
\texttt{1.  \$P(Xevu | \textbackslash{}neg Zuph, \textbackslash{}neg Rixq) = 0.15\$} \\
\texttt{2.  \$P(Xevu | \textbackslash{}neg Zuph, Rixq) = 0.50\$} \\
\texttt{3.  \$P(Xevu | Zuph, \textbackslash{}neg Rixq) = 0.32\$} \\
\texttt{4.  \$P(Xevu | Zuph, Rixq) = 0.85\$} \\
\\ 
\texttt{We are given \$P(Zuph) = 0.84\$, which implies \$P(\textbackslash{}neg} \\
            }
        }
    }
    & & \\ \\

    \theutterance \stepcounter{utterance}  
    & & & \multicolumn{2}{p{0.3\linewidth}}{
        \cellcolor[rgb]{0.9,0.9,0.9}{
            \makecell[{{p{\linewidth}}}]{
                \texttt{\tiny{[GM$|$GM]}}
                \texttt{Yes} \\
            }
        }
    }
    & & \\ \\

    \theutterance \stepcounter{utterance}  
    & & & \multicolumn{2}{p{0.3\linewidth}}{
        \cellcolor[rgb]{0.9,0.9,0.9}{
            \makecell[{{p{\linewidth}}}]{
                \texttt{\tiny{[GM$|$GM]}}
                \texttt{yes} \\
            }
        }
    }
    & & \\ \\

    \theutterance \stepcounter{utterance}  
    & & & \multicolumn{2}{p{0.3\linewidth}}{
        \cellcolor[rgb]{0.9,0.9,0.9}{
            \makecell[{{p{\linewidth}}}]{
                \texttt{\tiny{[GM$|$GM]}}
                \texttt{game\_result = WIN} \\
            }
        }
    }
    & & \\ \\

\end{supertabular}
}

\end{document}
